%% HAND-WRITTEN fragment -- NOT generated by maestro2tex (one transient record
%% decoded by hand, not a Maestro results database; there is no maestro2tex
%% config for any *_eis_* fragment).  Precedent: fig_eis_error.tex.
%%
%% Drop-in source: ~/chips/castalia/simruns/eisfigs2_20260826/build/
%% fig5_codelevel.tex, standalone preamble, in-figure annotation block and
%% in-figure note removed -- both are in the caption here.  The eight plotted
%% values are inlined rather than tabled: they are four bar tops and six points,
%% every one of them printed in the caption.  Full precision lives in
%% eisfigs2_20260826/data/fig5a_attribution.csv and fig5b_phase.csv.
%%
%% Source of every number: ~/chips/castalia/simruns/eisdeep_20260826/
%%   code_report/codeA_re_n20.json -- the headline record, 256 conversions,
%%     mux interleaved slot 0 (transimpedance output) <-> slot 6 (RE), 294 us,
%%     869 565.2 S/s, 64 conversions per excitation period -> f_exc =
%%     13 586.9565 Hz exactly; excitation from the real 12-bit ladder (code pair
%%     2032/2064); 11.000 kOhm load at N = 20; max |V_OUT - vcm| 0.132 V of the
%%     +-0.80 V window; extracted converter, 4451 device lines.
%%   code_report/pair_F0_F6.json -- the two-record scheme (codeF0_n20 +
%%     codeF6_n20), the "two records, one grid" point.
%%   code_report/deskew_test.txt, csv/codelevel.csv, sampling_prediction.txt.
%% AC analysis of the same load and gain code: 11 032.43 ohm / +0.3603 deg.
%%   Z_analog  +0.1262 % / -0.0127 deg   (20 ns strobe grid, boxcar-decimated,
%%                                        no sampling -- the transient path)
%%   Z_sampled +0.9812 % / -2.8949 deg   (same analog nets at the apertures)
%%   Z_code    +0.3431 % / -2.9383 deg   (the ten-bit words themselves)
%% Steps +0.126 / +0.855 / -0.638 pp.  Inter-channel gain mismatch +0.6359 %
%% (628.00 codes/V on the current channel, fundamental 92.79 LSB; 624.04 on the
%% potential channel, fundamental 7.65 LSB) predicts -0.632 pp against -0.638
%% measured -> 0.006 pp for quantisation, dither and aperture together.
%% Aperture fitted at nominal + 0 ns, 0.78 LSB rms; the three per-cycle code
%% phasors agree to 0.0000 %.  Half-conversion grid offset 180*f/f_s = 2.8125
%% deg; deskewed -0.082 deg (codes -0.126); two records on one grid +0.083 deg
%% measured; pure sine on the same grids 0.000 deg.  +-0.573 deg caps an
%% uncorrected interleaved read at 2768 Hz = 0.00318*f_s.
%%
%% THE PURE-SINE ROW IS COMPUTED HERE, NOT SIMULATED: an ideal excitation on the
%% analytic load with a single-pole Z_t fitted to the record's own zt_ac, point
%% sampled on the two interleaved grids and demodulated exactly as the extractor
%% does.  The same model driven with a two-code square returns +0.2727 % /
%% -2.9802 deg against the campaign's own sampling_prediction.txt value of
%% +0.2728 % / -2.9807 deg at 64 conversions per period -- four decimals, which
%% is what licenses the row.  Every other point in the figure is measured.
%%
%% Caveats that had to survive into the caption: both panels are referred to the
%% AC ANALYSIS OF THE SAME LOAD AND GAIN CODE, whose own error against the
%% analytic load is +0.295 % and +0.405 deg (against the analytic load the codes
%% read +0.639 % and -2.533 deg); ONE record, ONE die, ONE frequency, ONE
%% resistive load, no noise sources in the deck -- a deterministic attribution
%% and not a distribution, so there is nothing behind this figure to quote a
%% sigma from.  The record is demodulated over cycles 1-2; cycle 0 carries the
%% TIA's step response to the first DAC edge and the last cycle demodulates
%% 1.6 deg off in the boxcar path (eisdeep S10).  The complex-load code record
%% (codeC_rc_n0) saturated at 136 % of the compliance window and is excluded, so
%% this is a resistive-load result -- which is also why the attribution in (a) is
%% unambiguous.
%% Library cell names are kept out of the rendered prose.
\providecommand{\MaestroRoot}{include/analog/}
%% mtxA is the chapter's own palette slot A and carries the same value here.
%% mtxR and mtxN are the waterfall's diverging poles and a neutral total --
%% not chapter slots, so they carry names of their own; reusing a slot name
%% for a different colour is silently pre-empted by \providecolor.
\providecolor{mtxA}{HTML}{2A78D6}
\providecolor{mtxR}{HTML}{E34948}
\providecolor{mtxN}{HTML}{898781}
\begin{figure}[htbp]
  \centering
  {\pgfplotsset{compat=1.18}%
  \pgfplotsset{
    clax/.style={
      height=4.7cm,
      axis line style={line width=0.4pt, draw=black!45},
      tick style={line width=0.4pt, draw=black!45},
      tick align=outside,
      major grid style={line width=0.2pt, draw=black!14},
      label style={font=\small}, tick label style={font=\small},
      title style={font=\small},
    }}%
  \begin{tikzpicture}
    \begin{axis}[clax, width=0.40\linewidth,
      xmin=0.35, xmax=4.65, ymin=-0.30, ymax=1.20,
      ylabel={Error vs small signal~[points]},
      title={\textbf{(a) magnitude attribution}},
      xtick={1,2,3,4},
      xticklabels={{transient\\path},{sampled\\system},{gain\\mismatch},
                   {total}},
      xticklabel style={align=center, font=\scriptsize},
      ytick={0,0.25,0.5,0.75,1.0}, ymajorgrids,
    ]
      \filldraw[fill=mtxA, draw=white, line width=0.5pt]
        (axis cs:0.68,0) rectangle (axis cs:1.32,0.126203);
      \filldraw[fill=mtxA, draw=white, line width=0.5pt]
        (axis cs:1.68,0.126203) rectangle (axis cs:2.32,0.981157);
      \filldraw[fill=mtxR, draw=white, line width=0.5pt]
        (axis cs:2.68,0.981157) rectangle (axis cs:3.32,0.343062);
      \filldraw[fill=mtxN, draw=white, line width=0.5pt]
        (axis cs:3.68,0) rectangle (axis cs:4.32,0.343062);
      \draw[black!55, densely dotted, line width=0.5pt]
        (axis cs:1.32,0.126203) -- (axis cs:2.32,0.126203);
      \draw[black!55, densely dotted, line width=0.5pt]
        (axis cs:2.32,0.981157) -- (axis cs:3.32,0.981157);
      \draw[black!55, densely dotted, line width=0.5pt]
        (axis cs:3.32,0.343062) -- (axis cs:4.32,0.343062);
      \draw[black, line width=0.6pt] (axis cs:0.35,0) -- (axis cs:4.65,0);
      \node[font=\scriptsize, anchor=south] at (axis cs:1.0,0.161) {\(+0.126\)};
      \node[font=\scriptsize, anchor=south] at (axis cs:2.0,1.006) {\(+0.855\)};
      \node[font=\scriptsize, anchor=south] at (axis cs:3.0,1.006) {\(-0.638\)};
      \node[font=\scriptsize, anchor=south] at (axis cs:4.0,0.378) {\(+0.343\)};
    \end{axis}
  \end{tikzpicture}\hfill
  \begin{tikzpicture}
    \begin{axis}[clax, width=0.50\linewidth,
      xmin=-4.55, xmax=1.35, ymin=0.4, ymax=7.9,
      xlabel={Phase vs small signal~[\si{\degree}]},
      title={\textbf{(b) the phase term is the excitation}},
      xtick={-4,-3,-2,-1,0,1},
      ytick={1,2,3,4,5,6},
      yticklabels={{sinusoid control},{two records, one grid},
                   {deskewed},{no sampling},
                   {sampled},{codes}},
      yticklabel style={align=right, font=\scriptsize},
      xmajorgrids, ytick style={draw=none},
    ]
      \fill[black!7] (axis cs:-0.573,0.4) rectangle (axis cs:0.573,7.9);
      \draw[black!55, dashed, line width=0.6pt]
        (axis cs:-2.8125,0.4) -- (axis cs:-2.8125,7.9);
      \node[font=\scriptsize, anchor=north east, align=right]
        at (axis cs:-2.72,7.8) {\(-180f/f_s\)};
      \node[font=\scriptsize, anchor=north west] at (axis cs:-0.50,7.8)
        {\(\pm\SI{0.573}{\degree}\) gate};
      \addplot[only marks, mark=*, mark size=2.0pt, mtxA,
               mark options={fill=mtxA, draw=white, line width=0.4pt}]
        coordinates {(-2.9383,6) (-2.8949,5) (-0.0127,4) (0.0828,2)};
      \addplot[only marks, mark=o, mark size=2.0pt, black, line width=0.6pt]
        coordinates {(-0.0824,3) (0.0000,1)};
      \node[font=\scriptsize, anchor=east] at (axis cs:-2.9383,6) {\(-2.938\)~};
      \node[font=\scriptsize, anchor=east] at (axis cs:-2.8949,5) {\(-2.895\)~};
      \node[font=\scriptsize, anchor=west] at (axis cs:-0.0127,4) {~\(-0.013\)};
      \node[font=\scriptsize, anchor=west] at (axis cs:-0.0824,3) {~\(-0.082\)};
      \node[font=\scriptsize, anchor=west] at (axis cs:0.0828,2) {~\(+0.083\)};
      \node[font=\scriptsize, anchor=west] at (axis cs:0.0000,1) {~\(+0.000\)};
    \end{axis}
  \end{tikzpicture}}
  \caption[{Impedance demodulated from converter codes: magnitude attribution
  and the phase mechanism.}]{Impedance demodulated from the ten-bit words the
  extracted converter produced, on one record of \num{256} conversions at
  \SI{869565}{\per\second} with the input multiplexer interleaving the
  transimpedance output and the reference electrode, \num{64} conversions per
  excitation period at \SI{13.587}{\kilo\hertz}, the excitation synthesised by
  the real \num{12}-bit ladder, an \SI{11.000}{\kilo\ohm} load at gain code 20
  and \SI{0.132}{\volt} of the \(\pm\SI{0.80}{\volt}\) compliance window.
  \textbf{(a)} The same record read three ways attributes the magnitude error
  completely: the transient path reproduces the small-signal analysis to
  \(+0.126\) points, reading those same nets at the conversion apertures adds
  \(+0.855\), and decoding the codes takes \(0.638\) back --- which the
  converter's measured inter-channel gain mismatch of \(+0.636\,\%\)
  (\num{628.00} codes per volt on the current channel against \num{624.04} on
  the potential channel, which swings over \(12\times\) less of the transfer
  curve) predicts as \(-0.632\), leaving \num{0.006} points for quantisation,
  dither and aperture together. \textbf{(b)} The \SI{-2.94}{\degree} is the
  excitation and not the converter: interleaving samples the two channels half a
  conversion apart, \(180\,f/f_s = \SI{2.813}{\degree}\) here, and a sinusoid on
  the same two grids returns \SI{0.000}{\degree} exactly, so the term is edge
  quantisation of the two-level excitation; filled marks are measured, open
  rings the corrected reading and the analytic control, and the band is the
  \(\pm\SI{0.573}{\degree}\) gate, which caps an uncorrected interleaved read at
  \SI{2.77}{\kilo\hertz}. Both panels are referred to the small-signal analysis
  of the same load and gain code, whose own error against the analytic load is
  \(+0.295\,\%\) and \SI{0.405}{\degree}. \textbf{This is one record on one die
  at one frequency into a resistive load}, with no noise sources in the deck ---
  a deterministic attribution with no distribution behind it, in which every
  residual is quantisation, linearity or sampling and none of it is noise.}
  \label{fig:eis-codelevel}
\end{figure}
